\documentclass[12pt]{scrartcl}
\usepackage{polski}
\usepackage[utf8]{inputenc}

\usepackage{amsmath}
\usepackage{amssymb}
\usepackage{makecell}
\usepackage[dvipsnames]{xcolor}

\usepackage{pgfplots}
	\pgfplotsset{compat=1.15}
\usepackage{xifthen}
\usepackage{enumitem}
\usepackage{titletoc,tocloft}
	\cftsetindents{section}{0in}{0.2in}
	\cftsetindents{subsection}{0in}{0.4in}
	\cftsetindents{subsubsection}{0in}{0.6in}
	\cftsetindents{paragraph}{0in}{0.8in}
\usepackage{hyperref}
\hypersetup{
	colorlinks,
	citecolor=black,
	filecolor=black,
	linkcolor=black,
	urlcolor=black
}

\usepackage{pgfkeys}
\pgfkeys{
	/splot/.is family, /splot,
	default/.style = 
	{xlabel = $k$, ylabel = $h$},
	xlabel/.estore in = \xlabel,
	ylabel/.estore in = \ylabel,
}

\setcounter{secnumdepth}{5}
\setcounter{tocdepth}{4}

\newcommand{\link}[2]{
	\hyperref[#1]{\textit{\underline{#2}}}
}
\newcommand{\code}[1]{\texttt{#1}}
\newcommand{\myparagraph}[1]{\paragraph{#1}\mbox{}\\}
\newcommand{\fancy}[1]{\mathcal{#1}}
\newcommand{\macierz}[1]{\begin{bmatrix}#1\end{bmatrix}}
\newcommand{\bgcolor}[2]{\colorbox{#1!30}{#2}}

\newcommand\mysplot[2][]{%
	\pgfkeys{#1}%
	...
}
\newcommand{\notka}[1]{
	\textcolor{gray}{\textbf{Notka:\\\indent} #1}
}
\setlength{\parskip}{0cm}

\newcommand{\splot}[3][]{
	\pgfkeys{/splot, default, #1}
	\begin{tikzpicture}
	\begin{axis}[
	xmin = -5,
	xmax = 10,
	ymin = -3,
	ymax = 3,
	axis lines = middle,
	axis line style = {->},
	xlabel = $\xlabel$,
	ylabel = $\ylabel$,
	xtick = data,
	ytick = data,
	x tick label style = {
		color=black,
	},
	title = {#2},
	title style = {font=\bfseries},
	]
	\addplot+[ybar, black, fill=black, mark=none, opacity=0.35] plot coordinates {#3};
	\end{axis}
	\end{tikzpicture}
}
\newcommand{\rownaniezmacierzax}[5]{
	\bgcolor{YellowGreen}{#1} &= \bgcolor{Yellow}{#2}\bgcolor{YellowGreen}{#3} + \bgcolor{Dandelion}{#4}\cdot\bgcolor{Turquoise}{#5}
}
\newcommand{\rownaniezmacierzaysimple}[3]{
	\bgcolor{Orchid}{#1} &= \bgcolor{Red}{#2}\bgcolor{YellowGreen}{#3}
}
\newcommand{\m}{\phantom{-}}

\title{Komputerowe Systemy Sterowania}
\subtitle{opracowanie}
\author{}
\date{}

\begin{document}
\maketitle
\begin{center}
	\textcolor{white}{
	$$\fancy{PISANE\ KRWIA\ I\ LZAMI,\ BYSCIE\ WY\ CIERPIEC\ NIE\ MUSIELI}$$
}
\end{center}
\newpage

\tableofcontents

\newpage

\section{Wstęp}
Opracowanie ma na celu uporządkować wiedzę z innych opracowań, która jest niezbędna do tego przedmiotu.\\ Opracowanie.pdf oraz Opracowanie2.pdf zawierają dodatkowe opisy niektórych zagadnień, więc warto tam poszukać, jeśli nie zrozumieliście czegoś z tego opracowania.


\section{Wzory}
	\subsection{Splot}
    	Układu $\fancy{LNP}$:
    		\begin{equation}\label{eq:splot}
    			y(n)=\sum_{k=0}^{n} u(k)h(n-k) = \sum_{k=0}^{n} u(n-k)h(k)=u(n)*h(n)
    		\end{equation}
    	Układu $\fancy{LN}$:
    		\begin{equation}
    			y(n)=\sum_{-\infty}^{\infty} u(k)h(n-k) = \sum_{-\infty}^{\infty} u(n-k)h(k)=u(n)*h(n)
    		\end{equation}

	\subsection{Równanie stanu}\label{eq:rownanie_stanu}
	\begin{equation}
		\begin{cases}
		x(n+1) &= Ax(n) + Bu(n) + B_{w}w(n)\\
		y(n) &= Cx(n) + Du(n) + D_{v}v(n)
		\end{cases}
	\end{equation}\\
	
	\textbf{Oznaczenia}:\\
	$A$ -- zestaw wektorów decydujących o przejściu do następnego stanu (macierz tranzycji $x$),
	\\$B$ -- macierz wejścia/sterowania,
	\\$C$ -- macierz obserwacji/ wyjścia,
	\\$D$ -- macierz równa 0 nie istniejąca w naturze wystepuje w układach 
	analogowych; w układach dyskretnych nie powinno występować, ale czasem 
	się zdarza, że istnieje ze względu na szumy,
	\\$B_w$, $D_u$ -- macierze oddziaływania zakłóceń,
	\\$x+1$ -- reprezentuje następny stan -- wyliczamy go na podstawie poprzedniego stanu,
	\\$w$, $v$ -- sygnały przypadkowe. 
	\\
\notka{
	Z pełnego równania jednak zadań nie ma z tego co wiem, poza podaniem wzoru, więc oto uproszczone równanie stanu:\\
}
	\label{eq:rownanie_stanu_uproszczone}
	\begin{equation}
		\begin{cases}
			x(n+1) &= Ax(n) + Bu(n)\\
			y(n) &= Cx(n) + Du(n)
		\end{cases}
	\end{equation}\\

\subsection{Równanie różnicowe}\label{eq:rownanie_roznicowe}
	\begin{equation}
		y(n) = \sum\limits_{i=0}^{N}a_iu(n-i)
	\end{equation}




\section{Definicje}
	\subsection{Niezmienniczość}\label{sec:niezmienniczosc}
	
	$\fancy{R}: u(n)=y(n)=0 $ dla $ n<0 $ (stan spoczynku)\\
	$\fancy{R} \in \fancy{N} \Leftrightarrow
	\forall_u\forall_n \Big( y(n)=\fancy{R}u(n) \Rightarrow \forall_{k>0} \big( y(n-k)=\fancy{R}u(n-k)\big) \Big)$\\
	
	Niezależnie od czasu, kiedy podamy sygnał, układ na wyjściu zawsze udzieli tej samej odpowiedzi.
	Niezmienniczość to niezmienność działania.\\\\
\notka{
	Czas nie ma wpływu na wynik, liczy się tylko wejście i odpowiedź układu.
	\\Aby zamodelować układ $LP$ (nie: $LNP$), można zastosować \code{srand(time(NULL))}.
}
	\subsection{Liniowość}\label{sec:liniowosc}
	$ \fancy{R} \in \fancy{L} \Leftrightarrow w.l. $ (warunek liniowy)
	
	
	\[ w.l.
		\begin{cases}
		 	$$\fancy{R} \alpha u = \alpha \fancy{R}u\\
			\fancy{R} [ u_1 + u_2 ] = \fancy{R}u_1 + \fancy{R}u_2$$
		\end{cases}
	\]
	
	\[ w.l.
		\begin{cases}
			$$ \fancy{R}[ \alpha u_1 + \beta u_2 ] = \alpha \fancy{R}u_1 + \beta \fancy{R} u_2 $$
		\end{cases}
	\]
	\begin{center}
		$\forall_{u, u_1, u_2} \in \fancy{D}$, $\alpha, \beta \in \fancy{R}$
	\end{center}

	
\notka{
	Można podać pierwszą albo drugą wersję. Wersja druga jest po prostu połączeniem dwóch równań w jedno, obie są równoważne.\\
	Wejściowy przeskalowany sygnał da przeskalowane wyjście.\\
	Sygnały można składać.
}
	
\subsection{Przyczynowość}\label{sec:przyczynowosc}
	
	$ \fancy{R} \in \fancy{P} \Leftrightarrow \fancy{R}u_1(n) = \fancy{R}u_2(n) $\\
	$ \forall_{n \le k} $, $ \forall u_1(n), u_2(n) \in \fancy{D} $:\\
	\indent\indent\indent $ u_1(n) = u_2(n) $ dla $ n \le k $\\
	\indent\indent\indent $ u_1(n) \ne u_2(n) $ dla $ n > k $\\\\
\notka{
	Dopóki $u_1$ i $u_2$ są takie same, wyjście będzie takie samo. $u_1$ oraz $u_2$ są takie same do momentu $k$. Później niż w momencie $k$ sygnały wejściowe ($u_1$, $u_2$) się różnią, więc i sygnały wyjściowe ($y_1$, $y_2$) też się różnią.
}
	
\subsection{Stabilność}\label{sec:stabilnosc}
\indent\textbf{Definicja:}\\
$ \fancy{R} \in \fancy{S} \Leftrightarrow $\\
$ \forall_{u,n} \big(|u(n)|\le M < \infty \Rightarrow |y(n)| < \infty \big) $\\

\textbf{Warunek:}\\
$ \fancy{R} \in \fancy{L}, \fancy{N}, \fancy{P} \Rightarrow $\\
$ \fancy{R} \in \fancy{S} \Leftrightarrow l_1 \triangleq \sum\limits_{k=0}^{\infty}|h(k)|< \infty $\\

\textbf{Dowód na stabilność splotu:}\label{eq:stabilnoscsplotu}\\
$ |y(n)|=|h*u| \le \sum\limits_{k=0}^{\infty}|h(k)||u(n-k)|\le M l_1=M $\\

\newpage
\section{Zadania praktyczne}
\subsection{Splot z interpretacją graficzną}\label{sec:splot}
Obliczyć próbki odpowiedzi układu LNP o skończonej odpowiedzi impulsowej w postaci \{2, 1, 1, -2\} na wejście w postaci sekwencji \{1, 1, -1\}.

\textbf{Dane:}

$u(n) = \{1, 1, -1\}$

$h(n) = \{2, 1, 1, -2\}$
\\

\begin{tabular}{cc}
	\makecell{
		\splot[xlabel=k, ylabel=u]{ u(n) -- wejście }{ (0,1) (1,1) (2,-1) }
	}
	&
	\makecell{
		\splot{ h(n) -- odpowiedź impulsowa }{ (0,2) (1,1) (2,1) (3,-2) }
	}
\end{tabular}

% first output row
\begin{tabular}{cc}
	\makecell{
		\splot{ n=0, h(-k) }{ (-3,-2) (-2,1) (-1,1) (0,2) }
	}
	&
	\makecell{
		\splot{ n=1, h(1-k) }{ (-2,-2) (-1,1) (0,1) (1,2) }
	}
\end{tabular}

% second output row
\begin{tabular}{cc}
	\makecell{
		\splot{ n=2, h(2-k) }{ (-1,-2) (0,1) (1,1) (2,2) }
	}
	&
	\makecell{
		\splot{ n=3, h(3-k) }{ (0,-2) (1,1) (2,1) (3,2) }
	}
\end{tabular}

% third output row
\begin{tabular}{cc}
	\makecell{
		\splot{ n=4, h(4-k) }{ (1,-2) (2,1) (3,1) (4,2) }
	}
	&
	\makecell{
		\splot{ n=5, h(5-k) }{ (2,-2) (3,1) (4,1) (5,2) }
	}
\end{tabular}

% fourth output row
\begin{tabular}{cc}
	\makecell{
		\splot{ n=6, h(6-k) }{ (3,-2) (4,1) (5,1) (6,2) }
	}
	&
	\makecell{
		\splot{ n=7, h(7-k) }{ (4,-2) (5,1) (6,1) (7,2) }
	}
\end{tabular}

\textbf{Podsumowując obliczenia:}
\begin{align*}
y(0) &= u(0) \cdot h(0) \\	
	&= 1 \cdot 2=2 \\
y(1) &= u(0) \cdot h(0) + u(1) \cdot h(1) \\	
	&= 1 \cdot 1 + 1 \cdot 2=1 + 2 = 3\\
y(2) &= u(0) \cdot h(0) + u(1) \cdot h(1) + u(2) \cdot h(2) \\	
	&= 1\cdot1 
	+ 1\cdot1 
	+ (-1)\cdot2 
	= 1 +1 -2 = 0 \\
y(3) &= u(0) \cdot h(0) + u(1) \cdot h(1) + u(2) \cdot h(2) + u(3) \cdot h(3) \\	
	&= 1\cdot(-2)
	+ 1\cdot1
	+ (-1)\cdot1
	+ 0\cdot2
	= -2 +1 -1 +0 = -2 \\
y(4) &= u(0) \cdot h(0) + u(1) \cdot h(1) + u(2) \cdot h(2) + u(3) \cdot h(3) + u(4) \cdot h(4)\\
	&=
	  1		\cdot 0
	+ 1		\cdot (-2)
	+ (-1)	\cdot 1
	+ 0		\cdot 1
	+ 0		\cdot 2
	= 0 -2 -1 +0 +0 +0 = -3\\
y(5) &= u(0) \cdot h(0) + u(1) \cdot h(1) + u(2) \cdot h(2) + u(3) \cdot h(3) + u(4) \cdot h(4) + u(5) \cdot h(5)\\
	&=
	1		\cdot 0
	+ 1		\cdot 0
	+ (-1)	\cdot (-2)
	+ 0		\cdot 1
	+ 0		\cdot 1
	+ 0		\cdot 2
	= 0 +0 +2 +0 +0 +0 = 2\\
y(6) &= u(0) \cdot h(0) + u(1) \cdot h(1) + u(2) \cdot h(2) + u(3) \cdot h(3) +
	\\ &+ u(4) \cdot h(4) + u(5) \cdot h(5) + u(6) \cdot h(6)\\
	&=
	1		\cdot 0
	+ 1		\cdot 0
	+ (-1)	\cdot 0
	+ 0		\cdot (-2)
	+ 0		\cdot 1
	+ 0		\cdot 1
	+ 0		\cdot 2 = \\
	&= 0 +0 +0 +0 +0 +0 +0 = 0\\
y(7) &= u(0) \cdot h(0) + u(1) \cdot h(1) + u(2) \cdot h(2) + u(3) \cdot h(3) +
	\\ &+ u(4) \cdot h(4) + u(5) \cdot h(5) + u(6) \cdot h(6) + u(7) \cdot h(7)\\
	&=
	1		\cdot 0
	+ 1		\cdot 0
	+ (-1)	\cdot 0
	+ 0		\cdot 0
	+ 0		\cdot (-2)
	+ 0		\cdot 1
	+ 0		\cdot 1
	+ 0		\cdot 2 = \\
	&= 0 +0 +0 +0 +0 +0 +0 +0 = 0
\end{align*}
\notka{Wyniki można po prostu dać pod odpowiednimi wykresami, ale dałem pod spodem, żeby lepiej się sformatował dokument.}\\

\begin{tabular}{c c}
	\makecell[b]{
		\textbf{Wyniki:}\\
		y(0) = 2 \\
		y(1) = 3 \\
		y(2) = 0 \\
		y(3) = -2 \\
		y(4) = -3 \\
		y(5) = 2 \\
		y(6) = 0 \\
		y(7) = 0 \\\\\\\\\\
		$ $
	}
	&
	\makecell[b]{
		\splot[xlabel=n, ylabel=y]{ Wyniki w reprezentacji graficznej: }{
			(0,2) (1,3) (2,0) (3,-2) (4,-3) (5,2) (6,0) (7,0)
		}
	}
\end{tabular}

\subsection{Splot (odpowiedź układu LN)}\label{sec:splotln}

W razie czego jakby mi się już nie chciało tego robić:\\

Zadanie liczymy jak zwykły splot, ale bierzemy wszystkie próbki, jeżeli mamy dane. h(k) odwracamy jak zwykle i liczymy dla wszystkich próbek, jeżeli coś jest.

Innymi słowy zmienia się przedział, z którego bierzemy próbki:\\
$y(n)=\sum\limits_{-\infty}^{\infty} u(k)h(n-k) = \sum\limits_{-\infty}^{\infty} u(n-k)h(k)=u(n)*h(n)$\\
Czyli z przedziału $ <-\infty; \infty> $.

\subsection{Równanie układu (zadanie z macierzą)}\label{sec:macierz}
\def \vA {\macierz{\m0 & \m1\\ -1 & -1 }}

\textbf{Treść:} Wyznacz odpowiedź układu cyfrowego na 3 jedynki zakładając model stanowy SISO o wartościach początkowych $x(0)=\begin{bmatrix}1\\ 0\end{bmatrix}$ zadany macierzami $A = \vA$, $B = \begin{bmatrix}0 \\ 1 \end{bmatrix}$, $C = \begin{bmatrix}1 & 1 \end{bmatrix}$, D nie ma.\\

Jako że D nie ma, będziemy korzystać ze zmodyfikowanego \link{eq:rownanie_stanu_uproszczone}{wzoru uproszczonego:}

$$
	\begin{cases}
		\rownaniezmacierzax{$x(n+1)$}{$A$}{$x(n)$}{$B$}{$u(n)$}\\
		\rownaniezmacierzaysimple{$y(n)$}{$C$}{$x(n)$}
	\end{cases}
$$
\\

\textbf{Dane:}\\

$\bgcolor{Turquoise}{$u(n)$}=\{1, 1, 1\}$ -- dla kolejnych u są same zera\\
$\bgcolor{YellowGreen}{$x(0)$}=\begin{bmatrix}1\\ 0\end{bmatrix}$\\
$\bgcolor{Yellow}{$A$} = \vA$\\
$\bgcolor{Dandelion}{$B$} = \begin{bmatrix}0 \\ 1 \end{bmatrix}$\\
$\bgcolor{Red}{$C$} = \begin{bmatrix}1 & 1 \end{bmatrix}$\\

$n=0$\\
$$
\begin{cases}
	\rownaniezmacierzax{$x(1)$}{$A$}{$x(0)$}{$B$}{$u(n)$}\\
	\rownaniezmacierzaysimple{$y(0)$}{$C$}{$x(0)$}
\end{cases}
$$
$$
\begin{cases}
	\rownaniezmacierzax{$x(1)$}
	{$\vA$}							% A, nie zmienia się
	{$\macierz{1 \\ 0}$}			% x(n)
	{$\macierz{0 \\ 1}$}			% B
	{1}								% u(n)
	=\macierz{\m0 \\ -1} + \macierz{0 \\ 1}
	=\bgcolor{YellowGreen}{$\macierz{0 \\ 0}$}\\
	
	\rownaniezmacierzaysimple{$y(0)$}
	{$\macierz{1 & 1}$}				% C, nie zmienia się
	{$\macierz{1 \\ 0}$}			% x(n)
	= \bgcolor{Orchid}{$\macierz{1}$}
\end{cases}
$$\\

$n=1$\\
$$
\begin{cases}
	\rownaniezmacierzax{$x(2)$}
	{$\vA$}							% A, nie zmienia się
	{$\macierz{0 \\ 0}$}			% x(n)
	{$\macierz{0 \\ 1}$}			% B
	{1}								% u(n)
	=\macierz{0 \\ 0} + \macierz{0 \\ 1}
	=\bgcolor{YellowGreen}{$\macierz{0 \\ 1}$}\\
	
	\rownaniezmacierzaysimple{$y(1)$}
	{$\macierz{1 & 1}$}				% C, nie zmienia się
	{$\macierz{0 \\ 0}$}			% x(n)
	= \bgcolor{Orchid}{$\macierz{0}$}
\end{cases}
$$

$n=2$\\
$$
\begin{cases}
\rownaniezmacierzax{$x(3)$}
{$\vA$}							% A, nie zmienia się
{$\macierz{0 \\ 1}$}			% x(n)
{$\macierz{0 \\ 1}$}			% B
{1}								% u(n)
=\macierz{\m1 \\ -1} + \macierz{0 \\ 1}
=\bgcolor{YellowGreen}{$\macierz{1 \\ 0}$}\\

\rownaniezmacierzaysimple{$y(2)$}
{$\macierz{1 & 1}$}				% C, nie zmienia się
{$\macierz{0 \\ 1}$}			% x(n)
= \bgcolor{Orchid}{$\macierz{1}$}
\end{cases}
$$

$n=3$\\
$$
\begin{cases}
\rownaniezmacierzax{$x(4)$}
{$\vA$}							% A, nie zmienia się
{$\macierz{1 \\ 0}$}			% x(n)
{$\macierz{0 \\ 1}$}			% B
{0}								% u(n)
=\macierz{-1 \\ \m0} + \macierz{0 \\ 0}
=\bgcolor{YellowGreen}{$\macierz{\m0 \\ -1}$}\\

\rownaniezmacierzaysimple{$y(3)$}
{$\macierz{1 & 1}$}				% C, nie zmienia się
{$\macierz{1 \\ 0}$}			% x(n)
= \bgcolor{Orchid}{$\macierz{1}$}
\end{cases}
$$

$n=4$\\
$$
\begin{cases}
\rownaniezmacierzax{$x(5)$}
{$\vA$}							% A, nie zmienia się
{$\macierz{\m0 \\ -1}$}			% x(n)
{$\macierz{0 \\ 1}$}			% B
{0}								% u(n)
=\macierz{-1 \\ \m1} + \macierz{0 \\ 0}
=\bgcolor{YellowGreen}{$\macierz{-1 \\ \m1}$}\\

\rownaniezmacierzaysimple{$y(4)$}
{$\macierz{1 & 1}$}				% C, nie zmienia się
{$\macierz{\m0 \\ -1}$}			% x(n)
= \bgcolor{Orchid}{$\macierz{-1}$}
\end{cases}
$$

$n=5$\\
$$
\begin{cases}
\rownaniezmacierzax{$x(6)$}
{$\vA$}							% A, nie zmienia się
{$\macierz{-1 \\ 1}$}			% x(n)
{$\macierz{0 \\ 1}$}			% B
{0}								% u(n)
=\macierz{1 \\ 0} + \macierz{0 \\ 0}
=\bgcolor{YellowGreen}{$\macierz{1 \\ 0}$}\\

\rownaniezmacierzaysimple{$y(5)$}
{$\macierz{1 & 1}$}				% C, nie zmienia się
{$\macierz{-1 \\ 1}$}			% x(n)
= \bgcolor{Orchid}{$\macierz{0}$}
\end{cases}
$$

$n=6$\\
$$
\begin{cases}
\rownaniezmacierzax{$x(7)$}
{$\vA$}							% A, nie zmienia się
{$\macierz{1 \\ 0}$}			% x(n)
{$\macierz{0 \\ 1}$}			% B
{0}								% u(n)
=\macierz{\m0 \\ -1} + \macierz{0 \\ 0}
=\bgcolor{YellowGreen}{$\macierz{\m0 \\ -1}$}\\

\rownaniezmacierzaysimple{$y(6)$}
{$\macierz{1 & 1}$}				% C, nie zmienia się
{$\macierz{1 \\ 0}$}			% x(n)
= \bgcolor{Orchid}{$\macierz{1}$}
\end{cases}
$$

$n=7$\\
$$
\begin{cases}
\rownaniezmacierzax{$x(8)$}
{$\vA$}							% A, nie zmienia się
{$\macierz{\m0 \\ -1}$}			% x(n)
{$\macierz{0 \\ 1}$}			% B
{0}								% u(n)
=\macierz{-1 \\ 1} + \macierz{0 \\ 0}
=\bgcolor{YellowGreen}{$\macierz{-1 \\ 1}$}\\

\rownaniezmacierzaysimple{$y(7)$}
{$\macierz{1 & 1}$}				% C, nie zmienia się
{$\macierz{\m0 \\ -1}$}			% x(n)
= \bgcolor{Orchid}{$\macierz{-1}$}
\end{cases}
$$


\section{Treści zadań}
\subsection{2018/19}
\subsubsection{Termin 1}

\myparagraph{Podaj definicję niezmienniczości (grupa A)}
Patrz: \hyperref[sec:niezmienniczosc]{\textit{\underline{niezmienniczość}}}

\myparagraph{Podaj definicję stabilności (grupa A)}
Patrz: \hyperref[sec:stabilnosc]{\textit{\underline{stabilność}}}

\myparagraph{Podaj warunek stabilności (grupa B)}
Patrz: \hyperref[sec:przyczynowosc]{\textit{\underline{stabilność}}}

\myparagraph{Podaj definicję przyczynowości (grupa B)}
Patrz: \hyperref[sec:przyczynowosc]{\textit{\underline{przyczynowość}}}

\paragraph{Obliczyć próbki odpowiedzi układu LNP o skończonej odpowiedzi impulsowej}
\begin{itemize}[noitemsep,topsep=0pt]
		\item[a)] $h=\{2, 1, 1, -2\}$ na wejście w postaci sekwencji $u=\{1, 1, -1\}$ (grupa A)
		\item[b)] $h=\{-2, 1, 1, 2\}$ na wejście w postaci sekwencji $u=\{1, -1, 1\}$ (grupa B)
\end{itemize}
$ $\\
Patrz: \hyperref[sec:splot]{\textit{\underline{splot}}}

\subsubsection{Termin 2}

\myparagraph{Podaj definicję i warunek stabilności}
Patrz: \hyperref[sec:przyczynowosc]{\textit{\underline{stabilność}}}

\paragraph{Wymień i opisz modele matematyczne układów cyfrowych}
\begin{itemize}[noitemsep,topsep=0pt]
	\item[--] Splot -- patrz: \link{eq:splot}{wzór na splot}
	\item[--] Równanie stanu -- patrz: \link{eq:rownanie_stanu}{wzór na równanie stanu}
	\item[--] Równanie różnicowe -- patrz: \link{eq:rownanie_roznicowe}{wzór na równanie różnicowe}
\end{itemize}

\myparagraph{Wyznacz odpowiedź układu cyfrowego na 3 jedynki zakładając model stanowy SISO...}
Wyznacz odpowiedź układu cyfrowego na 3 jedynki zakładając model stanowy SISO o wartościach początkowych $\macierz{1 \\ 0}$ zadany macierzami $A$ = $\macierz{\m0\ \m1\\ -1 -1}$, $B = \macierz{0 \\ 1}$, $C = \macierz{1 & 1}$, $D$ nie ma.\\

\noindent Patrz: \link{sec:macierz}{zadanie z macierzą}

\subsubsection{Termin 3}
\myparagraph{Rozwiązać i napisać wzór końcowy SISO...}
Rozwiązać i napisać wzór końcowy SISO. Rozpisać x(n) i y(n) dla (chyba) 7 kolejnych wartości. Podane $A$, $B$, $C$, i jakieś udziwnione $u(n)$, założyć że $x(0) = 0$.\\

\noindent Patrz: \link{sec:macierz}{zadanie z macierzą}

\myparagraph{Splot układu LN (układ NIE był przyczynowy)}
Patrz: \hyperref[sec:splotln]{\textit{\underline{splot LN}}}

\subsubsection{Termin 4}
\myparagraph{Podaj definicję przyczynowości}

\myparagraph{Podaj definicję i warunek stabilności}

\myparagraph{Splot układu LNP}
u = $\{1, 1, 1\}$\\
h = $\{1, -1, 1\}$\\

\noindent Patrz: \link{sec:splot}{zadanie ze splotem}

\myparagraph{Zadanie z macierzą}
A = $\macierz{0 & 1\\ 1 & \frac{1}{2}}$\\
B = $\macierz{0 \\ 1}$\\
C = $\macierz{0 & 1}$\\
u = $\{1, 1, 1\}$\\
$x(0)$ = $\macierz{0 \\ 0}$\\

\noindent Patrz: \link{sec:macierz}{zadanie z macierzą}\\


\myparagraph{Uzupełnij tabelę}\\
\begin{tabular}{|c|c|c|}
\hline
    \textbf{--} & \textbf{ST} & \textbf{REK} \\
 \hline
    \textbf{S3} & & \\
\hline
    \textbf{S4} & & \\
\hline
\end{tabular}\\\\
ST -- czy system jest stabilny\\
REK -- czy system jest rekursywny\\
$S3$ -- system z zadania 3\\
$S4$ -- system z zadania 4\\

\noindent
\textbf{Rozwiązanie:}\\
\begin{tabular}{|c|c|c|}
\hline
    \textbf{--} & \textbf{ST} & \textbf{REK} \\
 \hline
    \textbf{S3} & T & N \\
\hline
    \textbf{S4} & N & T \\
\hline
\end{tabular}\\
\\
S3 jest stabilny, patrz: \link{eq:stabilnoscsplotu}{dowód na stabilność splotu}\\
S4 jest niestabilny (bo rozciąga się w czasie w nieskończoność)
\end{document} 
 
